\centerline{\bf \Huge Считаем уголки}

\begin{flushright}
        \scriptsize{Страх нужно побеждать.\\Дьяволо}
\end{flushright}

\problem{1}
Докажите, что внешний угол треугольника равен сумме двух углов треугольника, не смежных с ним.

\problem{2}
Дан треугольник $A B C$. Внутри треугольника отметили такую точку $O$, что $A O=$ $B O=C O$. Величины углов $A O B, B O C$ и $C O A$ равны $96^{\circ}, 138^{\circ}, 126^{\circ}$. Чему равны углы треугольника $A B C$ ?

\problem{3}
В равнобедренном треугольнике $A B C$ с основанием $B C \angle A=36^{\circ}$. Проведена биссектриса $B K$. Докажите, что $B K=B C$.

\problem{4}
В треугольнике $D E F$ проведена медиана $D K$. Найдите углы треугольника, если $\angle K D E$ $=70^{\circ}, \angle D K F=140^{\circ}$.

\problem{5}
Один угол равнобедренного треугольника в 4 раза больше другого. Найдите углы треугольника. (Необходимо разобрать все случаи.)

\problem{6}
Точки $M$ и $N$ лежат на стороне $A C$ треугольника $A B C$, причём $\angle A B M=\angle C$ и $\angle C B N$ $=\angle A$. Докажите, что треугольник $B M N$ равнобедренный.